\begin{thebibliography}{99}
\setlength{\itemsep}{1.5pt}\small
\bibitem{r4} Rescio G, Manni A, Ciccarelli M, Papetti A, Caroppo A, Leone A. A Deep Learning-Based Platform for Workers’ Stress Detection Using Minimally Intrusive Multisensory Devices. Sensors 2024;24:947. https://doi.org/10.3390/s24030947.
\bibitem{r5} Donati M, Olivelli M, Giovannini R, Fanucci L. ECG-Based Stress Detection and Productivity Factors Monitoring: The Real-Time Production Factory System. Sensors 2023;23:5502. https://doi.org/10.3390/s23125502.
\bibitem{r61} M. Khayyat M, Munshi RM, Alabduallah B, Lamoudan T, Ghith E, Kim T, et al. An improved biometric stress monitoring solution for working employees using heart rate variability data and Capsule Network model. PLOS ONE 2024;19:e0310776. https://doi.org/10.1371/journal.pone.0310776.
\bibitem{r3} Sadeghi M, McDonald AD, Sasangohar F. Posttraumatic stress disorder hyperarousal event detection using smartwatch physiological and activity data. PLOS ONE 2022;17:e0267749. https://doi.org/10.1371/journal.pone.0267749.
\bibitem{r28} Whiston A, Igou ER, Fortune DG, Analog Devices Team, Semkovska M. Examining Stress and Residual Symptoms in Remitted and Partially Remitted Depression Using a Wearable Electrodermal Activity Device: A Pilot Study. IEEE J Transl Eng Health Med 2023;11:96--106. https://doi.org/10.1109/JTEHM.2022.3228483.
\bibitem{r115} Cajal D, De La Cámara C, Posadas-De Miguel M, Torrijos N, Nadal Ó, Blanco T, et al. Evaluation of Stress Response Using Smartphone PPG for Anxiety and Depression Monitoring. IEEE Access 2025;13:212689--703. https://doi.org/10.1109/ACCESS.2025.3643644.
\bibitem{r7} Rafi H, Benezeth Y, Song FY, Reynaud P, Arnoux E, Demonceaux C. Tree-Based Personalized Clustered Federated Learning: A Driver Stress Monitoring Through Physiological Data Case Study. IEEE Internet Things J 2025;12:1688--98. https://doi.org/10.1109/JIOT.2024.3462302.
\bibitem{r8} AlArnaout Z, Zaki C, Kotb Y, AlAkkoumi M, Mostafa N. Exploiting heart rate variability for driver drowsiness detection using wearable sensors and machine learning. Sci Rep 2025;15:24898. https://doi.org/10.1038/s41598-025-08582-2.
\bibitem{r45} Pavan K, Singh A, Pawar DS, Ganapathy N. Multimodal Wearable-Based Automated Driver Inattention State Assessment Using Multidevices and Novel Cross-Modal Attention Framework. IEEE Sens Lett 2025;9:1--4. https://doi.org/10.1109/LSENS.2025.3596610.
\bibitem{r30} Fernandez J, Martínez R, Innocenti B, López B. Contribution of EEG Signals for Students’ Stress Detection. IEEE Trans Affect Comput 2025;16:1235--46. https://doi.org/10.1109/TAFFC.2024.3503995.
\bibitem{r74} Chen Q, Lee BG. Deep Learning Models for Stress Analysis in University Students: A Sudoku-Based Study. Sensors 2023;23:6099. https://doi.org/10.3390/s23136099.
\bibitem{r99} Gedam S, Dutta S, Jha R. Analyzing mental stress in Indian students through advanced machine learning and wearable technologies. Sci Rep 2025;15:20610. https://doi.org/10.1038/s41598-025-06918-6.
\bibitem{r10} Tutunji R, Kogias N, Kapteijns B, Krentz M, Krause F, Vassena E, et al. Detecting Prolonged Stress in Real Life Using Wearable Biosensors and Ecological Momentary Assessments: Naturalistic Experimental Study. J Med Internet Res 2023;25:e39995. https://doi.org/10.2196/39995.
\bibitem{r13} Toshnazarov K, Lee U, Kim BH, Mishra V, Najarro LAC, Noh Y. SOSW : Stress Sensing With Off-the-Shelf Smartwatches in the Wild. IEEE Internet Things J 2024;11:21527--45. https://doi.org/10.1109/JIOT.2024.3375299.
\bibitem{r35} Presby DM, Jasinski SR, Capodilupo ER. Wearable derived cardiovascular responses to stressors in free-living conditions. PLOS ONE 2023;18:e0285332. https://doi.org/10.1371/journal.pone.0285332.
\bibitem{r12} Campanella S, Altaleb A, Belli A, Pierleoni P, Palma L. A Method for Stress Detection Using Empatica E4 Bracelet and Machine-Learning Techniques. Sensors 2023;23:3565. https://doi.org/10.3390/s23073565.
\bibitem{r19} Hongn A, Bosch F, Prado LE, Ferrández JM, Bonomini MP. Wearable Physiological Signals under Acute Stress and Exercise Conditions. Sci Data 2025;12:520. https://doi.org/10.1038/s41597-025-04845-9.
\bibitem{r25} Kumar S, Raj Chauhan A, Akhil, Kumar A, Yang G. Resp-BoostNet: Mental Stress Detection From Biomarkers Measurable by Smartwatches Using Boosting Neural Network Technique. IEEE Access 2024;12:149861--74. https://doi.org/10.1109/ACCESS.2024.3461588.
\bibitem{r43} Barki H, Nkenyereye L, Chung W-Y. Detection and Classification of Mental Stress Using In-Ear Plethysmography and a Vision Transformer. IEEE Sens J 2025;25:4015--27. https://doi.org/10.1109/JSEN.2024.3512595.
\bibitem{r31} Mai N-D, Chung W-Y. On-Chip Mental Stress Detection: Integrating a Wearable Behind-The-Ear EEG Device With Embedded Tiny Neural Network. IEEE J Biomed Health Inform 2025;29:1872--85. https://doi.org/10.1109/JBHI.2024.3519600.
\bibitem{r14} Mohammadi A, Fakharzadeh M, Baraeinejad B. An Integrated Human Stress Detection Sensor Using Supervised Algorithms. IEEE Sens J 2022;22:8216--23. https://doi.org/10.1109/JSEN.2022.3157795.
\bibitem{r108} Pei X, Ghandehari A, Chakoma S, Rajendran J, Tavares-Negrete JA, Esfandyarpour R. A quantitative, multimodal wearable bioelectronic device for comprehensive stress assessment and sub-classification. Nat Commun 2026;17:1150. https://doi.org/10.1038/s41467-025-67747-9.
\bibitem{r79} Stojchevska M, Steenwinckel B, Van Der Donckt J, De Brouwer M, Goris A, De Turck F, et al. Assessing the added value of context during stress detection from wearable data. BMC Med Inform Decis Mak 2022;22:268. https://doi.org/10.1186/s12911-022-02010-5.
\bibitem{r16} Xu J, Song C, Yue Z, Ding S. Facial Video-Based Non-Contact Stress Recognition Utilizing Multi-Task Learning With Peak Attention. IEEE J Biomed Health Inform 2024;28:5335--46. https://doi.org/10.1109/JBHI.2024.3412103.
\bibitem{r17} Rescio G, Manni A, Caroppo A, Ciccarelli M, Papetti A, Leone A. Ambient and wearable system for workers’ stress evaluation. Comput Ind 2023;148:103905. https://doi.org/10.1016/j.compind.2023.103905.
\bibitem{r118} Vos G, Trinh K, Sarnyai Z, Rahimi Azghadi M. Generalizable machine learning for stress monitoring from wearable devices: A systematic literature review. Int J Med Inform 2023;173:105026. https://doi.org/10.1016/j.ijmedinf.2023.105026.
\bibitem{r52} Zhu L, Spachos P, Ng PC, Yu Y, Wang Y, Plataniotis K, et al. Stress Detection Through Wrist-Based Electrodermal Activity Monitoring and Machine Learning. IEEE J Biomed Health Inform 2023;27:2155--65. https://doi.org/10.1109/JBHI.2023.3239305.
\bibitem{r72} Almadhor A, Sampedro GA, Abisado M, Abbas S, Kim Y-J, Khan MA, et al. Wrist-Based Electrodermal Activity Monitoring for Stress Detection Using Federated Learning. Sensors 2023;23:3984. https://doi.org/10.3390/s23083984.
\bibitem{r1} Alsahreef MS, Alfahmi M, Alharbi E, Jones JM, Jaber M, Alomainy A. IoMT-Based Proactive Approach for Detecting Stress in Physicians to Reduce Medical Errors. IEEE Access 2026;14:45934--47. https://doi.org/10.1109/ACCESS.2026.3674740.
\bibitem{r56} Anandan M, Palanisamy R. Wearable ECG Sensor-Based Cognitive Stress Detection During Physical Activity Using Multifractal Detrended Fluctuation Analysis. IEEE Sens Lett 2025;9:1--4. https://doi.org/10.1109/LSENS.2025.3605112.
\bibitem{r50} Ali MS, Motin MA, Kabir S. ReTeNet: A Residual Encoder and Transformer Encoders Network for Stress Monitoring From Wearable Device. IEEE Signal Process Lett 2025;32:3635--9. https://doi.org/10.1109/LSP.2025.3607767.
\bibitem{r55} Ali MS, Bipro SS, Motin MA, Kabir S, Sharma M, Chowdhury MEH. TEANet: A Transpose-Enhanced Autoencoder Network for Stress Monitoring on Resource-Constrained Wearable Device. IEEE Access 2026;14:9560--71. https://doi.org/10.1109/ACCESS.2025.3644400.
\bibitem{r47} Ali MS, Motin MA, Bipro SS, Kabir S, Syfullah MK, Kumar DK. Detection of Mental Stress From Photoplethysmography for Under-Resourced Healthcare Systems. IEEE Sens J 2026;26:6293--301. https://doi.org/10.1109/JSEN.2025.3649586.
\bibitem{r58} Rashid N, Mortlock T, Faruque MAA. Stress Detection Using Context-Aware Sensor Fusion From Wearable Devices. IEEE Internet Things J 2023;10:14114--27. https://doi.org/10.1109/JIOT.2023.3265768.
\bibitem{r26} Awada M, Becerik Gerber B, Lucas GM, Roll SC. Stress appraisal in the workplace and its associations with productivity and mood: Insights from a multimodal machine learning analysis. PLOS ONE 2024;19:e0296468. https://doi.org/10.1371/journal.pone.0296468.
\bibitem{r38} Rahman FN, Bindra PS, Nawar A, Crane HT, Sanchez-Perez JA, Berkebile JA, et al. A wearable system enabling acute stress monitoring and closed-loop mitigation through transcutaneous median nerve stimulation. Biosens Bioelectron 2026;301:118461. https://doi.org/10.1016/j.bios.2026.118461.
\bibitem{r91} Jui SJJ, Deo RC, Acharya R, Barua PD, Soar J, Devi A. Explainable stress detection system using hybrid CNN-Transformer and uncertainty quantification with copula GAN-based data. Biomed Signal Process Control 2026;120:109934. https://doi.org/10.1016/j.bspc.2026.109934.
\bibitem{r87} Modi N, Kumar Y, Mehta K, Chaplot N. Physiological signal-based mental stress detection using hybrid deep learning models. Discov Artif Intell 2025;5:166. https://doi.org/10.1007/s44163-025-00412-8.
\bibitem{r42} Li Y, Li K, Chen J, Wang S, Lu H, Wen D. Pilot Stress Detection Through Physiological Signals Using a Transformer-Based Deep Learning Model. IEEE Sens J 2023;23:11774--84. https://doi.org/10.1109/JSEN.2023.3247341.
\bibitem{r70} Ehrhart M, Resch B, Havas C, Niederseer D. A Conditional GAN for Generating Time Series Data for Stress Detection in Wearable Physiological Sensor Data. Sensors 2022;22:5969. https://doi.org/10.3390/s22165969.
\bibitem{r105} Gabrielli D, Prenkaj B, Velardi P. Seamless monitoring of stress levels leveraging a foundational model for time sequences. Artif Intell Med 2026;173:103336. https://doi.org/10.1016/j.artmed.2025.103336.
\bibitem{r34} Bello-Orgaz G, Menéndez HD, Quintana-Quiroga F, Díaz-Álvarez A. Smart devices and stress level detection: Balancing machine learning model complexity and device intrusiveness. Biomed Signal Process Control 2026;117:109683. https://doi.org/10.1016/j.bspc.2026.109683.
\bibitem{r54} Momeni N, Valdes AA, Rodrigues J, Sandi C, Atienza D. CAFS: Cost-Aware Features Selection Method for Multimodal Stress Monitoring on Wearable Devices. IEEE Trans Biomed Eng 2022;69:1072--84. https://doi.org/10.1109/TBME.2021.3113593.
\bibitem{r51} Khan HA, Nguyen TN, Shafiq G, Mirza J, Javed MA. A Secure Wearable Framework for Stress Detection in Patients Affected by Communicable Diseases. IEEE Sens J 2023;23:981--8. https://doi.org/10.1109/JSEN.2022.3204586.
\bibitem{r22} Moser MK, Ehrhart M, Resch B. An Explainable Deep Learning Approach for Stress Detection in Wearable Sensor Measurements. Sensors 2024;24:5085. https://doi.org/10.3390/s24165085.
\bibitem{r53} Shikha S, Sethia D, Indu S. Optimization of Wearable Biosensor Data for Stress Classification Using Machine Learning and Explainable AI. IEEE Access 2024;12:169310--27. https://doi.org/10.1109/ACCESS.2024.3463742.
\bibitem{r86} Rahman S, Udhayakumar R, Kaplan D, McCarthy B, Dawood T, Mellor N, et al. Photoplethysmography as a noninvasive surrogate for microneurography in measuring stress-induced sympathetic nervous activation -- A machine learning approach. Comput Biol Med 2025;185:109522. https://doi.org/10.1016/j.compbiomed.2024.109522.
\bibitem{r29} De Marchi B, Agovi E, Aliverti A. Stress Level Assessment by a Multi-Parametric Wearable Platform: Relevance of Different Physiological Signals. Inf Syst Front 2025;27:113--37. https://doi.org/10.1007/s10796-024-10550-6.
\bibitem{r73} Onim MSH, Thapliyal H, Rhodus EK. Utilizing Machine Learning for Context-Aware Digital Biomarker of Stress in Older Adults. Information 2024;15:274. https://doi.org/10.3390/info15050274.
\bibitem{r48} Moser MK, Resch B, Ehrhart M. An Individual-Oriented Algorithm for Stress Detection in Wearable Sensor Measurements. IEEE Sens J 2023;23:22845--56. https://doi.org/10.1109/JSEN.2023.3304422.
\bibitem{r76} Bin Heyat MB, Akhtar F, Abbas SJ, Al-Sarem M, Alqarafi A, Stalin A, et al. Wearable Flexible Electronics Based Cardiac Electrode for Researcher Mental Stress Detection System Using Machine Learning Models on Single Lead Electrocardiogram Signal. Biosensors 2022;12:427. https://doi.org/10.3390/bios12060427.
\bibitem{r41} Ribeiro G, Postolache O, Martín FF. A New Intelligent Approach for Automatic Stress Level Assessment Based on Multiple Physiological Parameters Monitoring. IEEE Trans Instrum Meas 2024;73:1--14. https://doi.org/10.1109/TIM.2023.3342218.
\bibitem{r24} Subathra P, Malarvizhi S, Ferents Koni Jiavana K, Patil S. A Wearable Electronic Band for Stress Understanding Using Machine Learning. IEEE Sens J 2025;25:38639--48. https://doi.org/10.1109/JSEN.2025.3590380.
\bibitem{r101} Vos G, Trinh K, Sarnyai Z, Rahimi Azghadi M. Ensemble machine learning model trained on a new synthesized dataset generalizes well for stress prediction using wearable devices. J Biomed Inform 2023;148:104556. https://doi.org/10.1016/j.jbi.2023.104556.
\bibitem{r21} El Arwadi L, Abu Daher L. CardioMind: Deep Learning Framework for Stress Detection Using IoT Devices. IEEE Access 2025;13:197027--42. https://doi.org/10.1109/ACCESS.2025.3633760.
\bibitem{r102} Bloomfield LSP, Fudolig MI, Kim J, Llorin J, Lovato JL, McGinnis EW, et al. Predicting stress in first-year college students using sleep data from wearable devices. PLOS Digit Health 2024;3:e0000473. https://doi.org/10.1371/journal.pdig.0000473.
\bibitem{r33} Li J, Washington P. A Comparison of Personalized and Generalized Approaches to Emotion Recognition Using Consumer Wearable Devices: Machine Learning Study. JMIR AI 2024;3:e52171. https://doi.org/10.2196/52171.
\bibitem{r75} Abdul Kader L, Al-Shargie F, Tariq U, Al-Nashash H. One-Channel Wearable Mental Stress State Monitoring System. Sensors 2024;24:5373. https://doi.org/10.3390/s24165373.
\bibitem{r37} Kim H, Song S, Cho BH, Jang DP. Deep learning-based stress detection for daily life use using single-channel EEG and GSR in a virtual reality interview paradigm. PLOS ONE 2024;19:e0305864. https://doi.org/10.1371/journal.pone.0305864.
\bibitem{r49} Laiti J, Liu Y, Dunne PJ, Byrne E, Zhu T. Real-World Classification of Student Stress and Fatigue Using Wearable PPG Recordings. IEEE Trans Affect Comput 2026;17:587--601. https://doi.org/10.1109/TAFFC.2025.3628467.
\bibitem{r109} Kim W, Kutsuzawa G, Maruyama M. Emotion recognition and forecasting from wearable data via cluster-guided attention with cross-species pretraining. Intell Syst Appl 2025;27:200560. https://doi.org/10.1016/j.iswa.2025.200560.
\end{thebibliography}